\documentclass[10pt,a4paper]{amsart}
\usepackage[margin=19mm]{geometry}
\usepackage{amsmath,amssymb,mathtools}
\usepackage[scr=rsfs,cal=euler]{mathalfa}
\usepackage{xfrac}
\usepackage[unicode,hidelinks,bookmarks=false]{hyperref}
\newcommand{\ie}{i.e.\ }
\newcommand{\cf}{cf.\ }
\newcommand{\resp}{respectively\ }
\newcommand{\Subsection}{\subsection}
\providecommand{\nobreakdash}{\nobreak\hbox{-}}
\setlength{\emergencystretch}{2em}
\multlinegap0pt

\theoremstyle{plain}
\newtheorem{theorem}{Theorem}[section]
\newtheorem{lemma}[theorem]{Lemma}
\newtheorem{claim}[theorem]{Claim}
\newtheorem{proposition}[theorem]{Proposition}
\newtheorem{corollary}[theorem]{Corollary}
\newtheorem{question}{Question}
\newtheorem*{problem}{Problem}
\newtheorem{conjecture}{Conjecture}
\theoremstyle{definition}
\newtheorem*{definition}{Definition}
\newtheorem{example}[theorem]{Example}
\theoremstyle{remark}
\newtheorem{remark}[theorem]{Remark}
\title[Sylow synchronization in finite groups: the good case]{Sylow synchronization in finite groups: the good case}
\author{Hong Yi Huang, Francesca Lisi, Aluna Rizzoli, Luca Sabatini}
\date{Source: arXiv:2609.28403v1 (CC BY 4.0)}
\begin{document}
\maketitle

\noindent\emph{Source and licence.} The delivered work is the article shown in the title above, version arXiv:2609.28403v1 (CC BY 4.0), distributed under the Creative Commons Attribution 4.0 International licence, as recorded in the arXiv licence metadata for this exact version and shown on the abstract page saved with this item. The full licence notice, which carries the licence URL and the address of that abstract page, is the bundled file \texttt{licenses/arxiv-2609.28403-CC-BY-4.0.txt}.
\par\smallskip

\noindent\textit{Editorial scope.} Counted unit: the article's Proposition (source label prop:translated) (source0.tex lines 262--271, source label \texttt{prop:translated}): ``Let be a finite nilpotent group and a completely reducible faithful -module. Let be the Sylow subgroups of , and suppose that for each there exists such that . Then for every there exists such that C ''. It is delivered together with the article's complete original proof (source0.tex lines 286--398, one proof environment printed later in the article, detached from the statement by the article's own arrangement, 113 lines), copied line for line from the article's own text. Required context retained uncounted: none beyond the counted statement and proof. Editorial formatting only: an AMS \texttt{amsart} preamble that restates the article's own macro definitions and its own theorem declarations, and the theorem environments the retained lines use; the article's own class file is neither shipped nor input; the retained environments are renumbered sequentially in order of appearance, whereas the article prints them with its own numbering; the article's equation numbering is not reproduced and every cross-reference inside the retained text resolves to the wrapper's numbering. No citation occurs inside any retained line, so the wrapper carries no bibliography; All retained bodies are exact copies of the bundled \texttt{sources/211585/source0.tex}, so no omission receipt applies. No asset-inclusion command occurs inside any retained span and the package ships no image asset.


\section{The counted unit: Proposition (source label prop:translated) and its complete original proof}

\begin{proposition} \label{prop:translated}
Let $G$ be a finite nilpotent group and $V$ a completely reducible
faithful $G$-module.
Let $P_1,\ldots,P_n$ be the Sylow subgroups of $G=P_1\times\cdots\times P_n$,
and suppose that for each $i = 1,\ldots,n$ there exists $v_i \in V$ such that $C_{P_i}(v_i) = 1$.
Then for every $t_1,\ldots,t_n \in V$ there exists
$v\in V$ such that 
$$ C_{P_i}(v + t_i) \> = \> 1 $$
for every $i=1,\ldots,n$.
\end{proposition}

\begin{proof}[Proof of Proposition~\ref{prop:translated}]
For each isotypic component $W$ of $V$,
let $v_{i,W},t_{i,W}\in W$ be the projections of $v_i$ and $t_i$ respectively.
We shall find $x_W \in W$ such that
\begin{equation} \label{eq:component-containment}
    C_{P_i}(x_W+t_{i,W}) \subseteq C_{P_i}(v_{i,W})
    \qquad (1\leqslant i\le n).
\end{equation}
For $v=\sum_W x_W$, we have
\[
    C_{P_i}(v+t_i)
    =\bigcap_W C_{P_i}\bigl(x_W+ t_{i,W} \bigr)
    \subseteq \bigcap_W C_{P_i}\bigl(v_{i,W}\bigr)
    =C_{P_i}(v_i)=1
\]
for every $i$, as desired.

Now let $F$ be the ground field and fix an isotypic component $W$ of $V$.
We suppress the subscript $W$, so let $v_1,\ldots,v_n,t_1,\ldots,t_n \in W$.
Write $W \cong S^{\oplus m}$, where $S$ is an irreducible $FG$-module,
and set $E=\mathrm{End}_{FG}(S)$.
By Schur's lemma and Wedderburn's theorem, $E$ is a finite field.
Its action gives $S$ an $E$-vector-space structure,
and the image of $FG$ on $S$ is $\mathrm{End}_E(S)$.
For each $i$, let $A_i \cong \mathrm{Mat}_{d_i}(E)$ be the $E$-algebra generated by the image of $P_i$ on $S$.
These algebras are semisimple, commute pairwise, and together generate $\mathrm{End}_E(S)$.
After omitting the $P_i$ acting trivially on $W$
and relabelling the remaining indices, we have
\[
    W\cong U_1\otimes_E\cdots\otimes_E U_k\otimes_E E^m,
\]
where $P_i$ acts only on $U_i$, absolutely irreducibly, and $\dim_E U_i=d_i$.

Let $q=|E|$.
The nontrivial image of $P_i$ on $U_i$ has a central element of order $p_i$,
which acts as a scalar by absolute irreducibility.
Hence $p_i\mid q-1$ for each $i\leqslant k$, and $k<q$ because these primes are distinct.

Choose bases in all tensor factors.
For $y \in W$, let $M_i(y)$ be its $i$th flattening:
the $d_i\times r_i$ matrix of the contraction map
\[
    \left(\bigotimes_{j\ne i}U_j\otimes_E E^m\right)^*
    \to U_i,
    \qquad
    r_i=m\prod_{j\ne i}d_j.
\]
An element of \(P_i\) fixes \(y\) if and only if it fixes every column of \(M_i(y)\). If \(M_i(y)\) has full row rank, then the elements of \(P_i\) fixing $y$ act trivially on \(U_i\) and hence on \(W\), so
\begin{equation}\label{eq:full-rank-centralizer}
    C_{P_i}(y) = C_{P_i}(W) .
\end{equation}
Since \(C_{P_i}(W)\subseteq C_{P_i}(v_i)\), \eqref{eq:full-rank-centralizer} gives \eqref{eq:component-containment} for $y = x +t_i$.
We shall use the elementary fact that a nonzero polynomial over $\mathbb{F}_q$ in $N$ variables, of degree less than $q$ in each variable, cannot vanish on all of $\mathbb{F}_q^N$.

\smallskip
\noindent\emph{Case 1: $d_i \leqslant r_i$ for every $i \leqslant k$.}
Let $X$ be a formal vector with indeterminates as the coordinates in the chosen tensor product basis.
For each $i$, choose a $d_i \times d_i$ minor of $M_i(X+t_i)$ and denote its determinant by $f_i(X)$.
Each $f_i$ is a nonzero polynomial of degree at most one in each coordinate.
Thus
\[
    \prod_{i=1}^k f_i(X)
\]
is nonzero and has degree at most $k<q$ in each coordinate.
Hence there exists $x\in W$ at which this product is nonzero,
and \eqref{eq:full-rank-centralizer} holds.

\smallskip
\noindent\emph{Case 2: $d_1>r_1$, after relabelling.}
There is at most one such index, since $r_i \geqslant d_j$ whenever $j\ne i$.
Put $R=U_2\otimes_E\cdots\otimes_E U_k\otimes_E E^m$, so that $W=U_1\otimes_E R$.
Choose a full-column-rank matrix
$A\in\mathrm{Mat}_{d_1\times r_1}(E)$ whose column space contains that of
$M_1(v_1)$. Let $B$ be an $r_1\times r_1$ matrix of independent
indeterminates, and define $x(B)$ by
\[
    M_1(x(B)+t_1)=AB.
\]
For every invertible specialization of $B$, the matrices $AB$ and
$A$ have the same column space. Hence
\begin{equation}\label{eq:first-factor}
    C_{P_1}(x(B)+t_1) \subseteq C_{P_1}(v_1).
\end{equation}

Let $y(B)$ be defined by $M_1(y(B))=AB$, and put $y=y(I)$.
Every $M_i(y)$ with $i>1$ has full row rank.
Indeed, otherwise some nonzero functional on $U_i$ would annihilate $y$ by contraction;
tensoring it with a nonzero functional on the remaining factors of $R$
would give a nonzero vector in the kernel
of the injective map $M_1(y) \colon R^* \to U_1$.
For each $i>1$, choose a $d_i\times d_i$ minor nonzero on $M_i(y)$,
and denote its determinant, evaluated on $M_i(x(B)+t_i)$, by $f_i(B)$.
Since
\[
    x(B)+t_i=y(B)+(t_i-t_1),
\]
the homogeneous part of degree $d_i$ of $f_i(B)$ is the corresponding minor of $M_i(y(B))$.

To bound the degree in an individual entry $B_{ab}$, write $A_a$ for column $a$ of $A$,
and let $e_b$ be the corresponding tensor product basis vector of $R$.
The coefficient of $B_{ab}$ in $y(B)$ is the pure tensor $A_a\otimes e_b$.
Its matrix in every flattening has rank at most one, so multilinearity of the determinant implies that each \(f_i(B)\) has degree at most one in \(B_{ab}\).
It follows that the product
\[
    \det(B) \prod_{i=2}^k f_i(B)
\]
is a nonzero polynomial of degree at most $k<q$ in each variable.
Choosing a specialization over $E$ at which it is nonzero,
\eqref{eq:first-factor} handles $P_1$,
while \eqref{eq:full-rank-centralizer} handles every $P_i$ with $i>1$.
Thus \eqref{eq:component-containment} holds in this case as well,
and the proof is complete.
\end{proof}

\end{document}
